\section{Modification of the bias term}
\label{sec:bias}

The BC-SVM of \citet{nakayama2017} subtracts the bias term $\kappa$ in
\eqref{eq:1.3}, which comes from the geometric representation. On the other hand,
Theorem~\ref{thm:1} shows that the deterministic part of the limiting Gram matrix
is the diagonal matrix $D$ whose components are $c_{i0}=\lim\tr(\Sigma_{i*})/\theta$,
that is, only the trace of the non-spiked part. The spiked part
$\sum_{s\le m_i}\lambda_{is}$ does not have a deterministic limit.

\begin{proposition}
\label{prop:1}
Assume \textup{(A1)} to \textup{(A4)} and \textup{(S)}. Then the deterministic
offset term of the discriminant function is given by
\[
\kappa_*=\frac{\tr(\Sigma_{1*})}{n_1}-\frac{\tr(\Sigma_{2*})}{n_2}
=\kappa-\Biggl(
\frac{1}{n_1}\sum_{s\le m_1}\lambda_{1s}
-\frac{1}{n_2}\sum_{s\le m_2}\lambda_{2s}
\Biggr).
\]
Hence the BC-SVM, which uses an estimator of $\tr(\Sigma_i)$ including the spikes,
overcorrects the bias by
$\sum_{s\le m_1}\lambda_{1s}/n_1-\sum_{s\le m_2}\lambda_{2s}/n_2$, which is of
order $\theta$ and is not negligible compared with the signal $\Delta=O(\theta)$.
\end{proposition}

\begin{proof}[Outline of proof]
We expand the discriminant function for general $n_1$ and $n_2$ in the same way as
\eqref{eq:7.3}. Since $w=\sum_p\alpha_p y_p u_p$, the deterministic term of the
discriminant function comes from the diagonal part $D$ of Theorem~\ref{thm:1}
through the Karush--Kuhn--Tucker conditions, and it has the form
$\sum_p\alpha_p^* y_p\cdot(\text{diagonal of }D)$. In the existing theory,
$\theta_i$'s are equalized within each class and this gives the form $\kappa$ in
\eqref{eq:1.3}. Under the spiked model, the diagonal term is
$\tr(\Sigma_{i*})/\theta$ from~\eqref{eq:3.2} and not $\tr(\Sigma_i)/\theta$, so
that the replacement occurs. The difference is
$\tr(\Sigma_i)-\tr(\Sigma_{i*})=\sum_{s\le m_i}\lambda_{is}$ by definition.
\end{proof}

\begin{remark}
\label{rem:6}
Spiked eigenvalue structures are commonly observed in actual high-dimensional data
such as gene expression data. Proposition~\ref{prop:1} suggests that, when the
BC-SVM is applied to such data, the bias correction works too strongly and the
separating hyperplane is shifted excessively towards the minority class. We can
expect an improvement by replacing the estimator of $\tr(\Sigma_i)$ with
$\tr(\Sigma_i)-\sum_{s\le\hat m_i}\tilde\lambda_{is}$, where $\tilde\lambda_{is}$
denotes the estimator of the spiked eigenvalue given by the NR methodology.
\end{remark}
